%% parts/p03/qe-categoricity-games.tex --- P3 quantifier elimination,
%% categoricity, back-and-forth, and Ehrenfeucht--Fraisse games.
%%
%% Migration accounting for source/formal-logic-reference.tex 7150--7493:
%% every definition, method, result, warning, and worked example is retained.
%% The ACF categoricity example is promoted to a cited result (S-011).
%% Exercise qe-categoricity-01 is absorbed by thm:p03-dlo-countable-categorical;
%% exercise 02 is retired after its QE-to-inclusion argument is instantiated
%% in the DLO corollary and its both-models hypothesis is retained explicitly;
%% exercise 03 is retired as a redundant practice problem after the general
%% back-and-forth method and two applications have been retained.
%% Repairs R-015, R-016, and R-030 are preserved in annex:p3.qe.1--.5.
%% Apparatus: A-040, A-041, A-122.  Depth: D-P4-001 and stub D-P4-006.

\chapter[QE, Categoricity, and Games]{Quantifier Elimination, Categoricity, and Games}
\label{sec:p03-qe-categoricity-games}

\section{Quantifier elimination and model completeness}
\label{sec:p03-quantifier-elimination}

\begin{hbdefinition}{def:p03-quantifier-elimination}{Quantifier elimination}
Let (T) be a theory in a signature \(\Lang\).  The theory \emph{eliminates
quantifiers} if, for every \(\Lang\)-formula \(\varphi(\bar x)\), there is a
quantifier-free \(\Lang\)-formula \(\psi(\bar x)\) such that
\[
  T\entails\forall\bar x\,
  \bigl(\varphi(\bar x)\leftrightarrow\psi(\bar x)\bigr).
\]
Thus the equivalence is relative to the models of \(T\), not to all
\(\Lang\)-structures.
\index{quantifier elimination}
\end{hbdefinition}

\begin{hbmethod}{def:p03-one-quantifier-reduction}{One-quantifier reduction}
It suffices to eliminate the displayed quantifier from every formula
\(\exists y\,\theta(y,\bar x)\) with quantifier-free \(\theta\).  Put Boolean
combinations into a convenient normal form, distribute the existential over
finite disjunctions, replace universal quantifiers by negated existentials,
rename bound variables, and induct outward through the formula.  The method
reduces a global elimination theorem to an explicit analysis of one
quantifier; it is not itself an eliminating algorithm until that analysis is
given for the theory at hand.
\end{hbmethod}

\begin{hbtheorem}{thm:p03-dlo-qe}{cited=Marker2002}{Quantifier elimination for dense linear orders}
Let \(\mathsf{DLO}\) be the first-order theory of nonempty dense linear orders
without endpoints in the signature \(\{<\}\).  Then \(\mathsf{DLO}\)
eliminates quantifiers.
\end{hbtheorem}

For one conjunction of order literals, terms involving the eliminated
variable supply lower bounds, upper bounds, equalities, and disequalities.  An
equality permits substitution.  With no equality, density and the absence of
endpoints reduce existence to finite compatibility among the bounds and the
forbidden parameter values.  The cited proof verifies every literal case and
then applies the one-quantifier method \cite{Marker2002}.

\begin{hbexample}{def:p03-dlo-one-quantifier}{A one-quantifier DLO witness}
Modulo \(\mathsf{DLO}\),
\[
  \exists y\,(x<y\land y<z)\quad\leftrightarrow\quad x<z.
\]
Transitivity gives the forward implication, and density gives the reverse
implication.
\end{hbexample}

\begin{hbdefinition}{def:p03-theory-completeness}{Three completeness notions}
A theory \(T\) is \emph{deductively complete} if for every sentence \(\sigma\)
in its signature either \(T\derives\sigma\) or \(T\derives\neg\sigma\).  It is
\emph{semantically complete} if either \(T\entails\sigma\) or
\(T\entails\neg\sigma\) for every such \(\sigma\).  Soundness and semantic
completeness of the fixed first-order calculus identify these two properties.

The theory \(T\) is \emph{model complete} if every embedding between models
of \(T\) is elementary.  This third property concerns formulas with parameters
along arbitrary embeddings; it is neither of the preceding decision
properties.
\index{model complete theory}\index{deductively complete theory}
\index{semantically complete theory}
\end{hbdefinition}

\begin{hblemma}{thm:p03-model-completeness-reduction}{proved=annex:p3.qe.1}{Reduction of model completeness to inclusions}
Let \(T\) be a theory such that \(\mathcal M\preccurlyeq\mathcal N\) whenever
\(\mathcal M\subseteq\mathcal N\) are both models of \(T\).  Then every
embedding between models of \(T\) is elementary, so \(T\) is model complete.
\end{hblemma}

The conclusion quantifies over arbitrary embeddings.  Factoring an embedding
through its image is therefore load-bearing; this is the repaired reduction
R-015, not a restriction of model completeness to literal inclusions.

\begin{hbcorollary}{thm:p03-dlo-complete-model-complete}{sketched=annex:p3.qe.2}{DLO completeness and model completeness}
The theory \(\mathsf{DLO}\) is deductively complete (equivalently,
semantically complete) and model complete.
\end{hbcorollary}

\begin{hbtheorem}{thm:p03-acfp-qe}{cited=Marker2002}{\(\mathsf{ACF}_p\) quantifier elimination}
For each fixed characteristic \(p\), where \(p=0\) or \(p\) is prime, the
theory \(\mathsf{ACF}_p\) of algebraically closed fields of characteristic
\(p\) eliminates quantifiers in the ring language
\(\{0,1,+,-,\mathbin{\cdot}\}\).
\end{hbtheorem}

The characteristic is part of the theory: the union customarily called
``ACF'' without a fixed characteristic is not complete.  Real closed fields,
by contrast, have literal elimination in the \emph{ordered}-ring language;
Presburger arithmetic needs congruence or divisibility predicates before an
elimination theorem is literal.  These language boundaries are made explicit
in the two catalogs below \cite{Marker2002,Presburger1929}.

\begin{hbdefinition}{def:p03-definable-skolem-functions}{Definable Skolem functions}
A theory \(T\) has \emph{definable Skolem functions} if, for every formula
\(\varphi(\bar x,y)\), some formula \(\theta(\bar x,y)\) satisfies
\[
 T\entails\forall\bar x\Bigl(
   \exists y\,\varphi(\bar x,y)\to
   \exists!y\,[\varphi(\bar x,y)\land\theta(\bar x,y)]\Bigr).
\]
This is a property of a theory: \(\theta\) uniformly chooses a definable
witness on every nonempty fiber.  It is separated here because neither
quantifier elimination nor model completeness alone supplies such a choice;
the next catalog uses the distinction.
\index{definable Skolem functions}
\end{hbdefinition}

For the last complexity cell of Table~\ref{tab:p03-qe-catalog},
\(\mathrm{STA}(s,t,a)\) denotes the class accepted by alternating Turing
machines using space \(O(s(n))\), time \(O(t(n))\), and at most \(O(a(n))\)
alternations; an asterisk leaves that resource unrestricted.  This
metalinguistic abbreviation is needed to report Kozen's exact Boolean-algebra
classification and is not a synonym for \(\mathsf{PSPACE}\).

\begin{landscape}
\begingroup\tiny\setlength{\tabcolsep}{3pt}
\begin{longtable}{@{}p{1.15in}p{1.05in}p{.62in}p{.65in}p{.95in}p{2.64in}@{}}
\caption{First-order theories master table (apparatus A-040).  ``Complete''
means deductively complete in the displayed language; conditional negative
entries expose their consistency assumptions.}
\label{tab:p03-first-order-theories}\\
\toprule
Theory and exact language & Complete? & Decidable? & Literal QE? & Categoricity summary & Exact scope or caveat \\
\midrule\endfirsthead
\toprule
Theory and exact language & Complete? & Decidable? & Literal QE? & Categoricity summary & Exact scope or caveat \\
\midrule\endhead
\(\mathsf{DLO}\) in \(\{<\}\) & Yes & Yes & Yes & \(\aleph_0\)-categorical & Not categorical in any uncountable cardinal. \\
\addlinespace
\(\mathsf{ACF}_p\), fixed \(p=0\) or prime, in \(\{0,1,+,-,\mathbin{\cdot}\}\) & Yes & Yes & Yes & Uncountably categorical & Categorical in every uncountable cardinal; not \(\aleph_0\)-categorical. \\
\addlinespace
\(\mathsf{RCF}\) in \(\{0,1,+,-,\mathbin{\cdot},<\}\) & Yes & Yes & Yes & No infinite categoricity & The order symbol is essential to the literal QE entry. \\
\addlinespace
\(\operatorname{Th}(\N;0,1,+,<)\) (Presburger arithmetic) & Yes & Yes & No & No infinite categoricity & QE holds after adding all congruence, equivalently divisibility, predicates. \\
\addlinespace
\(\operatorname{Th}(\N;S)\) in \(\{S\}\) & Yes & Yes & No & Uncountably categorical & Not \(\aleph_0\)-categorical; naming the unique initial point yields QE. \\
\addlinespace
Atomless Boolean algebras in \(\{0,1,\wedge,\vee,\neg\}\) & Yes & Yes & Yes & \(\aleph_0\)-categorical & Not categorical in uncountable cardinals. \\
\addlinespace
\(\mathsf{PA}\) in \(\{0,S,+,\times\}\) & No, if consistent & No, if consistent & No, if consistent & No infinite categoricity, if consistent & Effective axiomatizability and arithmetic strength activate incompleteness. \\
\addlinespace
True arithmetic \(\operatorname{Th}(\N;0,1,+,\times,<)\) & Yes, by definition & No; not c.e. & No & No infinite categoricity & It has nonstandard elementary equivalents; completeness is not categoricity. \\
\addlinespace
\(\mathsf{ZFC}\) in \(\{\in\}\) & No, if consistent & No, if consistent & No, if consistent & No infinite categoricity, conditional on consistency & The row presupposes that the consistency/model condition is met. \\
\bottomrule
\end{longtable}
\endgroup
\end{landscape}

The DLO, fixed-characteristic ACF, and RCF classifications are standard
consequences of their elimination theorems; the arithmetic negative rows use
the effective and incompleteness boundaries developed elsewhere in the
handbook.  Sources for the table include
  \cite{Marker2002,Steinitz1910,Hodges1993,Presburger1929,Shoenfield1967,Godel1931,
FerranteRackoff1979,Kozen1980}.  No row infers categoricity merely from
completeness.

\begin{landscape}
\begingroup\tiny\setlength{\tabcolsep}{2pt}
\begin{longtable}{@{}p{.65in}p{1.10in}p{1.20in}p{.60in}p{.70in}p{.80in}p{2.10in}@{}}
\caption{Quantifier-elimination catalog (apparatus A-041).  An em dash in the
definable-choice column records that this entry deliberately makes no claim;
it does not mean failure.}
\label{tab:p03-qe-catalog}\\
\toprule
Theory & Base language and literal QE & Eliminating expansion or witness & Model complete? & Definable Skolem functions? & Complete / decidable? & Calibrated complexity statement \\
\midrule\endfirsthead
\toprule
Theory & Base language and literal QE & Eliminating expansion or witness & Model complete? & Definable Skolem functions? & Complete / decidable? & Calibrated complexity statement \\
\midrule\endhead
\(\mathsf{DLO}\) & \(\{<\}\): yes & Finite lower/upper-bound order types & Yes & No: automorphisms fixing \(x\) prevent a definable choice above \(x\) & Yes / yes & Validity is \(\mathsf{PSPACE}\)-complete \cite{Kozen2006}. \\
\addlinespace
\(\mathsf{ACF}_p\) & Ring language, fixed \(p\): yes & Algebraic-geometric elimination; resultants witness the procedure & Yes & No: conjugate roots cannot be uniformly selected & Yes / yes & Grigor'ev's bound is polynomial in \(L^{n^{2a+1}}\), parameterized by formula length \(L\), variables \(n\), and alternations \(a\); no unqualified polynomial-time claim \cite{Grigoriev1987}. \\
\addlinespace
\(\mathsf{RCF}\) & Ordered-ring language: yes; pure-ring language: no literal claim & Sign conditions and Tarski--Seidenberg projection & Yes & Yes: semialgebraic definable choice & Yes / yes & Some QE outputs require doubly exponential time and space; this is not an exact decision-class classification \cite{DavenportHeintz1988}. \\
\addlinespace
Presburger & \(\{0,1,+,<\}\): no & Add congruence/divisibility predicates; Cooper-style elimination & Yes & --- & Yes / yes & Fischer--Rabin give a doubly exponential lower bound in formula length in their computation model \cite{FischerRabin1974}. \\
\addlinespace
\(\operatorname{Th}(\N;S)\) & \(\{S\}\): no & Add the initial-point constant \(0\); predecessor existence becomes \(x\ne0\) & No in the base language; yes after naming \(0\) & --- & Yes / yes & The one-successor theory in \(\{S\}\) is \(\mathsf{PSPACE}\)-complete \cite{FerranteRackoff1979}. \\
\addlinespace
Atomless Boolean algebras & \(\{0,1,\wedge,\vee,\neg\}\): yes & Split finite Boolean combinations into their Boolean atoms & Yes & --- & Yes / yes & Complete under log-space reductions for \(\bigcup_{c<\omega}\mathrm{STA}(*,2^{cn},n)\), including the atomless theory; do not relabel this class \(\mathsf{PSPACE}\) \cite{Kozen1980}. \\
\bottomrule
\end{longtable}
\endgroup
\end{landscape}

The shift embedding \(n\mapsto n+1\) maps \((\N,S)\) onto a substructure
whose first element has no predecessor internally but does have one in the
ambient structure; hence the bare successor theory is not model complete.
Naming \(0\) blocks that embedding, and quantifier elimination in
\(\{0,S\}\) gives model completeness.  The three em dashes in the choice
column keep that independent property open rather than manufacturing it from
QE.

\section{Categoricity and back-and-forth}
\label{sec:p03-categoricity-back-and-forth}

\begin{hbdefinition}{def:p03-categoricity}{Categoricity in a cardinal}
Let \(T\) be a theory and let \(\kappa\) be an infinite cardinal.  The theory
\(T\) is \emph{\(\kappa\)-categorical} if it has a model of cardinality
\(\kappa\) and any two models of \(T\) of cardinality \(\kappa\) are
isomorphic.
\index{categoricity}
\end{hbdefinition}

\begin{hbwarning}{def:p03-categoricity-cardinal}{Categoricity is cardinal-specific}
Categoricity in one cardinal is not categoricity in every cardinal.  The
L\"owenheim--Skolem theorems produce models in several infinite cardinalities;
they do not preclude uniqueness separately inside one cardinal.  For example,
the pure-equality theory asserting at least \(n\) elements for every
\(n\geq1\) is categorical in every infinite cardinal.  Transfer results such
as Morley's theorem require their own hypotheses.
\end{hbwarning}

\begin{hbdefinition}{def:p03-finite-partial-isomorphism}{Finite partial isomorphism}
A \emph{finite partial isomorphism} from \(\mathcal M\) to \(\mathcal N\) is a
finite injective map \(f:A\to N\), with \(A\subseteq M\), that preserves and
reflects every atomic formula on tuples from \(A\).
\index{partial isomorphism!finite}
\end{hbdefinition}

\begin{hbmethod}{def:p03-back-and-forth}{Back-and-forth}
For countable structures, construct an isomorphism as the union of an
increasing sequence of finite partial isomorphisms.  A ``forth'' step puts the
next enumerated element of the first structure into the domain; a ``back'' step
puts the next element of the second into the range.  The name denotes this
alternating construction, not a new map or relation.  Unlike the one-sided
extension property called homogeneity, the two directions force totality and
surjectivity together.
\index{back-and-forth method}
\end{hbmethod}

\begin{hbtheorem}{thm:p03-back-and-forth-test}{proved=annex:p3.qe.3}{Back-and-forth test}
Let \(\mathcal M\) and \(\mathcal N\) be countable structures.  Suppose
\(\mathcal F\) is a nonempty set of finite partial isomorphisms from
\(\mathcal M\) to \(\mathcal N\) such that, for every \(f\in\mathcal F\):
\begin{enumerate}
  \item for every \(a\in M\), some \(g\in\mathcal F\) extends \(f\) and has
        \(a\) in its domain; and
  \item for every \(b\in N\), some \(g\in\mathcal F\) extends \(f\) and has
        \(b\) in its range.
\end{enumerate}
Then, for \emph{every} starting map \(f_0\in\mathcal F\), there is an
isomorphism \(h:\mathcal M\to\mathcal N\) with \(f_0\subseteq h\), extension
taken as inclusion of sets of ordered pairs.  In particular,
\(\mathcal M\cong\mathcal N\).
\end{hbtheorem}

The universal extension clause in the conclusion is the R-016 repair.  It is
stronger than the bare assertion that some isomorphism exists and is required
by the next corollary.

\begin{hbcorollary}{thm:p03-countable-saturated-automorphism}{proved=annex:p3.qe.4}{Finite elementary self-maps of countable saturated structures}
Assume the Axiom of Choice in the ambient set theory.  If \(\mathcal M\) is
countable and \(\aleph_0\)-saturated, every finite elementary map from
\(\mathcal M\) to itself extends to an automorphism of \(\mathcal M\).
\end{hbcorollary}

\begin{hbtheorem}{thm:p03-dlo-countable-categorical}{proved=annex:p3.qe.5}{Countable categoricity of dense linear order}
Any two countable dense linear orders without endpoints are isomorphic.
Consequently, \(\mathsf{DLO}\) is \(\aleph_0\)-categorical.
\end{hbtheorem}

\begin{hbtheorem}{thm:p03-acfp-uncountable-categorical}{cited=Steinitz1910}{Uncountable categoricity of algebraically closed fields}
Fix \(p=0\) or a prime.  For every uncountable cardinal \(\kappa\), the theory
\(\mathsf{ACF}_p\) is \(\kappa\)-categorical.  It is not
\(\aleph_0\)-categorical.
\end{hbtheorem}

By the Steinitz classification, an algebraically closed field of fixed
characteristic and uncountable cardinality \(\kappa\) has transcendence degree
\(\kappa\) over its prime field, so equal-cardinality transcendence bases give
isomorphic algebraic closures.  In the countable cardinal, distinct finite or
countable transcendence degrees give nonisomorphic models
\cite{Steinitz1910,Marker2002}.

\begin{hbtheorem}{thm:p03-ryll-nardzewski}{cited=RyllNardzewski1959}{Ryll--Nardzewski theorem}
Let \(T\) be a deductively complete theory in a countable signature with an
infinite model.  The following are equivalent:
\begin{enumerate}
  \item \(T\) is \(\aleph_0\)-categorical;
  \item for every \(n\geq1\), the type space \(S_n(T)\) is finite;
  \item for every \(n\geq1\), every complete \(n\)-type over
        \(\varnothing\) is isolated; and
  \item for every \(n\geq1\), only finitely many formulas with free variables
        among \(x_1,\ldots,x_n\) exist up to equivalence modulo \(T\).
\end{enumerate}
\end{hbtheorem}

In this setting the unique countable model is prime.  It is also
\(\aleph_0\)-saturated.  The theorem's independent formulations and that
consequence are treated in
\cite{RyllNardzewski1959,Engeler1959,Svenonius1959,Hodges1993,Marker2002}.

\begin{hbtheorem}{thm:p03-morley-categoricity}{cited=Marker2002}{Morley's categoricity theorem}
Let \(T\) be a deductively complete theory in a countable signature.  If
\(T\) is categorical in one uncountable cardinal, then it is categorical in
every uncountable cardinal.
\end{hbtheorem}

\begin{hbwarning}{def:p03-completeness-versus-categoricity}{Theory completeness versus categoricity}
A deductively complete theory may have many nonisomorphic models of the same
cardinality.  Completeness says that the models agree on sentences;
\(\kappa\)-categoricity says that the models of the specified size are
isomorphic.
\end{hbwarning}

\hbstub{D-P4-006: a statement-level entry for Vaught's two-cardinal theorem,
including its exact two-cardinal hypotheses and a calibrated method sketch,
is deferred because this Phase-4 row is priority could; new-content row
N-a-005 remains open.}

\section{Ehrenfeucht--Fra\"iss\'e games}
\label{sec:p03-ef-games}

\begin{hbdefinition}{def:p03-quantifier-rank}{Quantifier rank and bounded agreement}
The \emph{quantifier rank} \(\operatorname{qr}(\varphi)\) is defined
recursively: atomic formulas have rank \(0\); negation preserves rank; a
binary Boolean connective takes the maximum rank of its arguments; and
\(\exists x\,\varphi\) and \(\forall x\,\varphi\) have rank
\(\operatorname{qr}(\varphi)+1\).  Two structures \emph{agree through rank
\(k\)} if they satisfy the same first-order sentences of quantifier rank at
most \(k\).  This metalinguistic relation is needed to state the finite-round
game theorem and is weaker than elementary equivalence, which quantifies over
sentences of every finite rank.
\index{quantifier rank}
\end{hbdefinition}

\begin{hbdefinition}{def:p03-ef-game}{Ehrenfeucht--Fra\"iss\'e game}
Fix a finite relational signature \(\tau\), two \(\tau\)-structures
\(\mathcal A,\mathcal B\), and \(k\in\N\).  In each of \(k\) rounds,
Spoiler chooses an element in either structure and Duplicator responds with an
element of the other.  Record every pair as \(a_i\mapsto b_i\), regardless of
which structure Spoiler chose from.  Duplicator wins exactly when the resulting
pairing is a finite partial isomorphism: it respects repetitions/equality and
preserves and reflects every relation symbol on selected tuples.  The finite
relational vocabulary is fixed; function symbols require a separate
relationalization or closure formulation.
\index{Ehrenfeucht--Fraisse game@Ehrenfeucht--Fra\"iss\'e game}
\end{hbdefinition}

\begin{hbtheorem}{thm:p03-ef}{cited=Ehrenfeucht1961}{Ehrenfeucht--Fra\"iss\'e theorem}
For structures \(\mathcal A\) and \(\mathcal B\) in the same fixed finite
relational signature, Duplicator has a winning strategy in the \(k\)-round
Ehrenfeucht--Fra\"iss\'e game if and only if \(\mathcal A\) and
\(\mathcal B\) agree through rank \(k\).
\end{hbtheorem}

The theorem converts syntactic rank into a combinatorial invariant; it does
not turn one fixed finite game into full elementary equivalence.  The game,
rank induction, and its finite-model applications are developed in
\cite{Ehrenfeucht1961,Libkin2004}.

\begin{hbcorollary}{thm:p03-ef-inexpressibility}{cited=Libkin2004}{Game criterion}
Let \(\mathcal C\) be an isomorphism-closed class of finite structures in one
fixed finite relational signature.  If, for every \(k\), there are finite
\(\mathcal A_k\in\mathcal C\) and \(\mathcal B_k\notin\mathcal C\) on which
Duplicator wins the \(k\)-round game, then no first-order sentence defines
\(\mathcal C\) on finite structures.
\end{hbcorollary}

\begin{figure}[htbp]
\centering
\begin{tikzpicture}[
  x=1.35cm,y=1.25cm,
  point/.style={circle,draw,minimum size=6.5mm,inner sep=1pt,font=\small},
  chosen/.style={circle,draw,double,minimum size=7.5mm,inner sep=1pt,font=\small},
  order/.style={-{Latex[length=2mm]}},
  pair/.style={dashed,thick}
]
  \node[point]  (a0) at (0,1.7) {\(a_0\)};
  \node[chosen] (a1) at (1,1.7) {\(a_1\)};
  \node[chosen] (a2) at (2,1.7) {\(a_2\)};
  \draw[order] (a0) -- (a1);
  \draw[order] (a1) -- (a2);
  \node[left=5mm of a0,font=\small] {\(\mathcal A\)};

  \node[point]  (b0) at (-.5,0) {\(b_0\)};
  \node[chosen] (b1) at (.5,0) {\(b_1\)};
  \node[point]  (b2) at (1.5,0) {\(b_2\)};
  \node[chosen] (b3) at (2.5,0) {\(b_3\)};
  \draw[order] (b0) -- (b1);
  \draw[order] (b1) -- (b2);
  \draw[order] (b2) -- (b3);
  \node[left=5mm of b0,font=\small] {\(\mathcal B\)};

  \draw[pair] (a1) -- node[left,font=\footnotesize] {round 1} (b1);
  \draw[pair] (a2) -- node[right,font=\footnotesize] {round 2} (b3);
\end{tikzpicture}
\caption{A two-round EF position on the three- and four-element linear orders
(apparatus A-122).  Spoiler first chose \(a_1\), answered by \(b_1\), and
then chose \(b_3\), answered by \(a_2\).  The selected pairing preserves
equality and order.  This board illustrates one play of a fixed-round strategy;
it is not the uniform family needed for the parity result.}
\label{fig:p03-ef-board}
\end{figure}

For the full two-round strategy in Figure~\ref{fig:p03-ef-board}, Duplicator
answers an endpoint by the corresponding endpoint and an interior first move
by an interior point.  On round two, the response repeats the paired element
when Spoiler repeats a move, and otherwise lies on the same side of the first
selected point.  With only two selected points this preserves their complete
order type.

\begin{hbexample}{def:p03-ef-parity}{Finite order parity is not first-order definable}
For \(k\geq1\), compare the linear orders of sizes \(2^k\) and \(2^k+1\).
With \(r\) rounds remaining, Duplicator maintains that the selected pairing is
order-preserving and every pair of corresponding gaps either has the same size
below \(2^r-1\), or both gaps have size at least \(2^r-1\).  Initially the
invariant holds.  A repeated selection receives its existing mate.  In equal
small gaps Duplicator copies the exact offset.
In two large gaps, if Spoiler leaves fewer than \(2^{r-1}-1\) points on one
side, Duplicator copies that smaller side; otherwise Duplicator chooses a point
leaving at least \(2^{r-1}-1\) points on both sides.  Thus the invariant
survives, and at \(r=0\) the selected map preserves equality and order.  The
\(k=0\) game is vacuous.  Hence every rank has an even/odd pair that agrees
through that rank, so Corollary~\ref{thm:p03-ef-inexpressibility} shows that no
first-order sentence defines even cardinality on finite linear orders.
\end{hbexample}

\begin{hbexample}{def:p03-ef-connectivity}{Finite graph connectivity is not first-order definable}
For each \(k\), let \(C_m\) denote the simple undirected cycle on \(m\)
vertices, and compare \(C_{2m}\) with the disjoint union \(C_m\sqcup C_m\)
for sufficiently large \(m\).  The cited EF argument matches the bounded
neighborhoods of selected vertices: with \(r\) rounds left it preserves all
distances below \(2^r\), while points outside those neighborhoods are answered
outside the corresponding neighborhoods.  For sufficiently large \(m\), the
responses remain available through \(k\) rounds.  One graph is connected and
the other is not, so connectivity is not first-order definable on finite
graphs.  The full invariant and response lemma are cited from
\cite{Libkin2004}; the preceding parity example supplies the fully expanded
strategy pattern used by this entry.
\end{hbexample}
